% !TeX program = xelatex
\documentclass[11pt,a4paper]{article}

\usepackage[a4paper,margin=0.62in]{geometry}
% Компилировать с XeLaTeX или LuaLaTeX.
\usepackage{fontspec}
\IfFontExistsTF{Lato}{\setmainfont{Lato}}{\setmainfont{Noto Sans}}
\usepackage{polyglossia}
\setdefaultlanguage{russian}
\usepackage{microtype}
\usepackage{tabularx}
\usepackage{array}
\usepackage{enumitem}
\usepackage{xcolor}
\usepackage{hyperref}
\usepackage{fancyhdr}
\usepackage{lastpage}

\definecolor{linkblue}{RGB}{20,33,110}
\hypersetup{
  colorlinks=true,
  urlcolor=linkblue,
  linkcolor=linkblue,
  citecolor=linkblue
}

\pagestyle{fancy}
\fancyhf{}
\renewcommand{\headrulewidth}{0pt}
\renewcommand{\footrulewidth}{0pt}
\cfoot{\small \thepage\ из \pageref{LastPage}}

\setlength{\parindent}{0pt}
\setlength{\parskip}{0pt}
\setlength{\tabcolsep}{0pt}
\setlist[itemize]{leftmargin=1.15em,itemsep=0.08em,topsep=0.15em,parsep=0pt,partopsep=0pt}

\newcolumntype{L}{>{\raggedright\arraybackslash}p{0.145\textwidth}}
\newcolumntype{R}{>{\raggedright\arraybackslash}p{0.835\textwidth}}

\newcommand{\cvrow}[2]{%
  \noindent
  \begin{minipage}[t]{0.145\textwidth}
    \vspace{0pt}%
    \raggedright
    \fontsize{10.5}{12.5}\selectfont
    #1
  \end{minipage}%
  \hfill
  \begin{minipage}[t]{0.835\textwidth}
    \vspace{0pt}%
    #2
  \end{minipage}\par
}

\newcommand{\entry}[2]{%
  \begin{tabularx}{\linewidth}{@{}X >{\raggedleft\arraybackslash}p{0.28\linewidth}@{}}
    #1 & #2
  \end{tabularx}
}

\newcommand{\pub}[2]{%
  #1 #2\par
  \vspace{0.28em}
}

\newcommand{\eduentry}[2]{%
  \noindent
  \begin{tabularx}{\linewidth}
    {@{}X@{\hspace{0.8em}}r@{}}
    #1 & \small #2
  \end{tabularx}\par
  \vspace{0.15em}
}

\newcommand{\edudetail}[1]{%
  \noindent\hspace*{1.15em}$\diamond$ #1\par
}

\newcommand{\skillcell}[2]{%
  \noindent
  \hangindent=1.15em
  \hangafter=1
  \makebox[1.15em][l]{$\diamond$}%
  \textbf{#1} #2%
}

\newcommand{\workentry}[2]{%
  \noindent
  \begin{tabularx}{\linewidth}
    {@{}X@{\hspace{0.8em}}r@{}}
    #1 & \small #2
  \end{tabularx}\par
  \vspace{0.15em}
}

\newcommand{\workproject}[1]{%
  \noindent
  \hangindent=2.30em
  \hangafter=1
  \hspace*{1.15em}\makebox[1.15em][l]{$\diamond$}#1\par
  \vspace{0.12em}
}

\newcommand{\workprojectdetail}[1]{%
  \noindent
  \hangindent=3.45em
  \hangafter=1
  \hspace*{2.30em}\makebox[1.15em][l]{--}#1\par
  \vspace{0.08em}
}

\begin{document}

% Шапка
{\fontsize{24}{28}\selectfont\bfseries Артур Ясеновец}\par
\vspace{0.08em}
{\large\bfseries Инженер-исследователь в области информационной безопасности}\par
\vspace{0.28em}
\href{mailto:archee.busy@gmail.com}{archee.busy@gmail.com}
\;\textbullet\;
\href{https://github.com/iasenovets}{GitHub}
\;\textbullet\;
\href{https://iasenovets.github.io/}{Портфолио}
\;\textbullet\;
\href{https://orcid.org/0009-0008-0008-3746}{ORCID}
\;\textbullet\;
\href{https://linkedin.com/in/iasenovets}{LinkedIn}
\par
\vspace{0.22em}
\hrule
\vspace{0.55em}

\cvrow{Профиль}{%
  Более 4 лет опыта в разработке программного обеспечения и инженерии
  информационной безопасности. Я разрабатываю на Go, Python, Rust и TypeScript,
  а также имею некоторый опыт работы с C++. Я переношу криптографические
  протоколы из научных статей в код, провожу анализ безопасности, корректности
  и утечек и разрабатываю системы для конфиденциального поиска, извлечения
  и обработки данных на основе гомоморфного шифрования.
}

\vspace{0.75em}

\cvrow{Навыки}{%
  \begin{tabularx}{\linewidth}
    {@{}>{\raggedright\arraybackslash}X
    @{\hspace{1.2em}}
    >{\raggedright\arraybackslash}X@{}}

    \skillcell{Программирование:}{%
      Go, Python, Rust, TypeScript, C++, SQL, Bash, LaTeX}
    &
    \skillcell{Криптографические библиотеки:}{%
      Lattigo, SEAL, TenSEAL, OpenFHE, fheanor}
    \\[0.20em]

    \skillcell{Технологии повышения конфиденциальности:}{%
      Приватное извлечение информации (PIR/CPIR),
      Симметричное шифрование с возможностью поиска (SSE),
      Приватный поиск шаблона}
    &
    \skillcell{Системы и инфраструктура:}{%
      Linux, Proxmox VE, VMware, Docker,
      GitHub Actions, GitLab CI/CD, Nginx, Ansible,
      Airflow, S3, Hyperledger Fabric}
    \\[0.20em]

    \skillcell{Гомоморфное шифрование:}{%
      Brakerski--Gentry--Vaikuntanathan (BGV),
      Brakerski--Fan--Vercauteren (BFV),
      Cheon--Kim--Kim--Song (CKKS)}
    &
    \skillcell{Инженерия безопасности:}{%
      YARA, Chainsaw, анализ вредоносного ПО/журналов,
      MITRE ATT\&CK}
    \\[0.20em]

    \skillcell{Постквантовая криптография:}{%
      ML-KEM, ML-DSA, SLH-DSA; LWE, RLWE}
    &
    \skillcell{Языки:}{%
      русский --- родной;
      английский --- \href{https://iasenovets.github.io/files/toefl_8cd5f4a2b8ec8f34da49ea0faeef43f1.pdf}{C1};
      китайский --- \href{https://admin.chinesetest.cn/queryScore.do?sid=c2tActR9Wb\%2704b861aefbdc5e9fd0b44afa717d08cbceb0720e582ce25f\%2724d7sWsQ5nVUQSwD}{HSK3}}
  \end{tabularx}
}

\vspace{0.75em}

\cvrow{Опыт работы}{%
  \workentry
    {\textbf{Магистрант-исследователь} @
    \href{https://www.cqupt.edu.cn/}{CQUPT}}
    {сент. 2024--наст. время}

  \workproject{Спроектировал и реализовал на Go с использованием Lattigo
  приватное извлечение информации на основе BGV для конфиденциальных запросов
  к мировому состоянию в Hyperledger Fabric; приватность запросов в модели
  HBC-узла основана на IND-CPA-безопасности BGV.}

  \workproject{Спроектировал архитектуру TWAPPA-FRL и реализовал с использованием
  TenSEAL приватную агрегацию обновлений модели DQN на основе CKKS; приватность
  обновлений клиентов в модели HBC-агрегатора основана на IND-CPA, при этом
  явно проанализированы утечки информации.}

  \workproject{Реализовал и проверил прототип приватного поиска шаблона в стиле
  FPS/RFPS, заменив преобразование сопоставления на основе DFT при приближённой
  арифметике CKKS преобразованием сопоставления на основе NTT при точной
  арифметике BGV/BFV; приватность шаблона в модели HBC-сервера основана
  на IND-CPA, при этом явно определены утечки протокола.}

  % приватность [запроса/обновления клиента/шаблона] на основе IND-CPA в модели HBC-[узла/агрегатора/сервера]

  \workproject{Восстановил модели предварительной обработки, построение запросов,
  серверные вычисления и развёртывание ответов для таких PIR-протоколов, как
  Piano, FrodoPIR/SimplePIR, XPIR, SealPIR и OnionPIR, и проанализировал
  компромиссы между объёмом коммуникации и производительностью.}
}
\vspace{0.50em}

\cvrow{}{%
  \workentry
    {\textbf{Инженер по информационной безопасности} @
    \href{https://global.ptsecurity.com/en/}{Positive Technologies}}
    {авг. 2024--февр. 2025}
  \workproject{Построил автоматизированный конвейер обработки вредоносного ПО
  с DAG в Airflow, который собирает образцы из MalwareBazaar, сохраняет
  артефакты в S3 и сканирует их с помощью YARA.}
  \workproject{Разработал на Python инструменты с покрытием модульными тестами
  для проверки путей Windows, сетевых конфигураций и политик контроля
  конечных устройств.}
  \workproject{Анализировал журналы EVTX, PowerShell и вредоносные макросы VBA;
  извлекал IoC и сопоставлял TTP с MITRE ATT\&CK для CVE-2020-5902
  и CVE-2023-38831.}
}

\vspace{0.50em}

\cvrow{}{%
  \workentry
    {\textbf{Инженер-программист} @ \href{https://academai.ru/}{AcademAI}}
    {сент. 2023--авг. 2024}
  \workproject{Разработал прикладной слой образовательной платформы на основе
  LLM, используемой преподавателями и студентами ИТ-направлений университетов:
  приложение на Next.js/TypeScript, сервисы на Python/Flask, слой данных
  PostgreSQL/Prisma и аутентификацию Auth.js.}
  \workproject{Интегрировал процессы генерации на основе LangChain с AsyncIO
  и Redis для асинхронной оркестрации LLM и управления промежуточным состоянием.}
  \workproject{Контейнеризировал приложение с помощью Docker, настроил Nginx,
  развернул его в облачной инфраструктуре VPS и создал процесс CI/CD
  в GitHub Actions.}
}

\cvrow{}{%
  \workentry
    {\textbf{Младший инженер-программист} @
    \href{https://zvekprogress.ru/}{ООО «ЗВЭК «Прогресс»»}}
    {май 2022--авг. 2023}

  \workproject{Разработал REST-бэкенд на C++17/Crow с MySQL, а также настольные
  и мобильные клиенты на Qt 5/QML для системы складского учёта на основе QR-кодов.}

  \workproject{Обеспечил поддержку приблизительно 1,000 записей и более
  20 одновременных пользователей; добавил покрытие модульными тестами, сборку
  с помощью CMake, разработку с использованием Git, отладку и развёртывание в Linux.}
}

\vspace{0.75em}

\cvrow{Образование}{%
  \raggedright

  \textbf{Чунцинский университет почты и телекоммуникаций},
  Чунцин, Китай\par
  \vspace{0.12em}

  \eduentry
    {$\diamond$ \textbf{Магистратура} по направлению «Сетевая и информационная безопасность»}
    {сент. 2024--июль 2027 \textit{(ожидается)}}

  \edudetail{Научный руководитель:
    \href{https://faculty.cqupt.edu.cn/tangfei/en/index.htm}{проф. Тан Фэй}}
  \edudetail{Средний балл: 3,98/4,00}

  \vspace{0.48em}

  \textbf{Ухтинский государственный технический университет},
  Ухта, Россия\par
  \vspace{0.12em}

  \eduentry
    {$\diamond$ \textbf{Бакалавриат} по направлению «Информационные системы и технологии»}
    {март 2022--июнь 2024}

  \edudetail{Средний балл: 4,49/5,00}

  \vspace{0.48em}

  \textbf{Санкт-Петербургский государственный университет телекоммуникаций},
  Санкт-Петербург, Россия\par
  \vspace{0.12em}

  \eduentry
    {$\diamond$ \textbf{Обучение в бакалавриате} по направлению «Информационные системы и технологии»}
    {сент. 2020--март 2022 \textit{(переведён)}}
}

\vspace{0.75em}

\cvrow{Публикации}{%
  \textbf{Опубликованные работы и препринты}\par\vspace{0.25em}

  \pub{\underbar{Artur Iasenovets}, Fei Tang, H. Zhu, P. Wang, and L. Liu.}
  {``\href{https://arxiv.org/abs/2511.02656}
  {Bringing Private Reads to Hyperledger Fabric via Private Information Retrieval},''
  arXiv:2511.02656, 2025. Рукопись находится на рецензировании.}

  \pub{H. Zhu, F. Tang, W. Qin, J. Shan, P. Wang, and \underbar{Artur Iasenovets}.}
  {``\href{https://doi.org/10.1007/s44443-025-00405-8}
  {EVMK-SSE: Efficient and Verifiable Multi-Keyword Symmetric Searchable Encryption},''
  \textit{Journal of King Saud University -- Computer and Information Sciences}, 2025.}

  \pub{\underbar{Artur Iasenovets} and Fei Tang.}
  {``\href{https://doi.org/10.60797/IRJ.2026.163.57}
  {MFCTIIF: Multi-Feed Cyber Threat Intelligence Integration Framework},''
  \textit{International Research Journal}, № 1 (163), 2026.}

  \pub{Piyumi M. S. Rajapaksha$^{*}$, \underbar{Artur Iasenovets$^{*}$},
  Yinxue Yi, and Fei Tang.}
  {``TWAPPA-FRL: Trust-Weighted Privacy-Preserving Aggregation for Federated
  Reinforcement Learning in Collaborative Intrusion Detection,''
  2026. Рукопись находится на рецензировании. $^{*}$Равный вклад.}
}

\vspace{0.75em}

\cvrow{Награды и\\достижения}{%
  \entry
    {\textbf{Стипендия правительства Китая (CSC) для обучения в магистратуре}}
    {2024--2027}
  Полная стипендия на весь период обучения по программе магистратуры в CQUPT.

  \vspace{0.45em}
  \entry
    {\textbf{Региональный демо-день Sber Student Accelerator}}
    {2023}
  Завершил программу Sber Student Accelerator в составе команды AcademAI
  и принял участие в региональном демо-дне Северо-Западного округа.

  \vspace{0.45em}
  \entry
    {\textbf{Победитель конкурса EdTech-проектов}}
    {2023}
  Представил образовательную платформу на основе LLM и победил в конкурсе,
  организованном Физтех-школой МФТИ, Фондом развития Физтех-школ
  и Startech.Base.

  \vspace{0.45em}
  \entry
    {\textbf{Phystech Gigachat Challenge от МФТИ и Сбера}}
    {2023}
  Принял участие в хакатоне Phystech Gigachat Challenge, разрабатывая
  образовательную платформу на основе LLM с использованием API GigaChat
  и Kandinsky в рамках 48-часового спринта.

  \vspace{0.45em}
  \entry{\textbf{Хакатон Kaspersky Cyberimmunity 2.0}}
  {2023}
  Принял участие в хакатоне Kaspersky Cyberimmunity осенью 2023 года
  и занял 11-е место среди более чем 100 команд, разработав решения
  для мониторинга трафика в целях безопасности низковысотного воздушного пространства.

  \vspace{0.45em}
  \entry
    {\textbf{Диплом I степени, «Севергеоэкотех»}}
    {2022}
  Занял первое место в секции «Современные информационные технологии»
  23-й Международной молодёжной научной конференции.

  \vspace{0.3em}
  \href{https://iasenovets.github.io/files/certificates_selected_en_public_ready_v2.pdf}
       {\small Подтверждающие сертификаты}
}

\end{document}
